\documentclass[10pt,twocolumn,letterpaper]{article}

%%%%%%%%% PAPER TYPE
% \usepackage{cvpr}              % CAMERA-READY
\usepackage[review]{cvpr}        % REVIEW  <-- current build
% \usepackage[pagenumbers]{cvpr} % arXiv version with page numbers

%% =====================================================================
%% PREAMBLE  (was preamble.tex)
%% Must be loaded before hyperref.
%% =====================================================================

\usepackage{multirow}
\usepackage{array}
\usepackage{makecell}
\usepackage{amsmath}
\usepackage{amssymb}
\usepackage{bm}
\usepackage{xspace}
\usepackage{microtype}   % tighter typesetting, roughly one sixth of a page
%% --- added for draft v1 -------------------------------------------------
\usepackage{xcolor}      % \color in \ph, \req, \ver (harmless if cvpr.sty already loads it)
\usepackage{graphicx}
\usepackage{booktabs}    % \toprule, \midrule, \bottomrule in the tables
\usepackage{pifont}      % \cmark / \xmark in Table 1
\newcommand{\cmark}{\ding{51}}
\newcommand{\xmark}{\ding{55}}
%% ------------------------------------------------------------------------

\newif\ifphfinal
\phfinalfalse          % <-- set to \phfinaltrue for the submission build

\ifphfinal
  \newcommand{\ph}[1]{\textbf{\fbox{UNFILLED: #1}}}
\else
  \newcommand{\ph}[1]{{\color{red}\texttt{[#1]}}}
\fi
\newcommand{\req}[1]{{\color{red}\textbf{[REQUIRED EXPERIMENT: #1]}}}
\newcommand{\ver}{{\color{red}$^{\dagger}$}}

\newcommand{\TODO}[1]{\textbf{\color{red}[TODO: #1]}}
\newcommand{\todo}[1]{{\color{red}#1}}

%% ---------------------------------------------------------------------
%% Notation
%% ---------------------------------------------------------------------
\newcommand{\csi}{\mathbf{H}}
\newcommand{\zwifi}{\mathbf{z}^{\mathrm{wifi}}}
\newcommand{\zvis}{\mathbf{z}^{\mathrm{vis}}}
\newcommand{\svis}{\hat{\mathbf{s}}^{\mathrm{vis}}}
\newcommand{\Lalign}{\mathcal{L}_{\mathrm{align}}}
\newcommand{\Ltask}{\mathcal{L}_{\mathrm{task}}}
\newcommand{\Lwave}{\mathcal{L}_{\mathrm{wave}}}
\newcommand{\brpm}{\,\mathrm{brpm}}
\newcommand{\bpm}{\,\mathrm{bpm}}

%% Method name. Change in one place.
\newcommand{\method}{RespWave\xspace}
%% Dataset name. Change in one place.
\newcommand{\dataset}{Meridian\xspace}

%% Compact table rules
\providecommand{\etal}{\textit{et al}.\xspace}
\providecommand{\ie}{\textit{i.e}.\xspace}
\providecommand{\eg}{\textit{e.g}.\xspace}

%% Tighter float spacing
\setlength{\textfloatsep}{9pt plus 2pt minus 2pt}
\setlength{\floatsep}{9pt plus 2pt minus 2pt}
\setlength{\intextsep}{9pt plus 2pt minus 2pt}

%% Figure placeholder until the real figure exists: \figph{width}{height}{what goes here}
\newcommand{\figph}[3]{\fbox{\parbox[c][#2][c]{#1}{\centering\small\color{red}#3}}}

\definecolor{cvprblue}{rgb}{0.21,0.49,0.74}
\usepackage[pagebackref,breaklinks,colorlinks,allcolors=cvprblue]{hyperref}

%%%%%%%%% PAPER ID
\def\paperID{*****}
\def\confName{CVPR}
\def\confYear{2027}

%%%%%%%%% TITLE
\title{\dataset: A Synchronized Contact and Contactless Physiological Dataset\\
under Emotion, Cognitive Load, Exercise, and Paced Breathing}

%%%%%%%%% AUTHORS
\author{Anonymous \confName~submission\\
Paper ID \paperID}

\begin{document}
\maketitle


%% =====================================================================
%% 0_abstract.tex
%% =====================================================================
\begin{abstract}
Contactless physiological sensing from cameras, thermal imagers, and Wi-Fi is usually trained and evaluated on
recordings of people sitting calmly, where heart rate stays in a narrow band and breathing is whatever the
participant happens to do. The conditions that make these signals useful are the ones those datasets leave out:
emotional arousal, cognitive load, exercise, and deliberate changes in breathing. We present \dataset, a dataset of
\ph{N} participants (\ph{H} hours) recorded by two lossless face cameras, a thermal camera, and Wi-Fi channel state
information (CSI), synchronized on a single host clock with contact references: chest ECG, finger
photoplethysmography, electrodermal activity, surface EMG, and audio. Every participant follows the same scripted
34-step, \ph{37}-minute protocol: six emotion film clips with per-clip self-reports, a cognitive stress task,
reading aloud, two exercises that raise heart rate, a recorded recovery, visually paced breathing at 6 and
20~breaths/min, and a guided relaxation. The protocol builds ground truth into the session: the pacer fixes the
breathing rate, the exercise sweeps heart rate across \ph{X--Y}\bpm, and the self-reports label every clip. We define
benchmarks for remote heart-rate estimation, contactless respiration from video, thermal and Wi-Fi, cross-modal
Wi-Fi sensing supervised by vision, and affect recognition, and we report baselines that show
\ph{one-sentence headline result, \eg rPPG error rises by X\bpm after exercise}. Data, code, and the recording
software will be released.
\end{abstract}


%% =====================================================================
%% 1_intro.tex
%% =====================================================================
\section{Introduction}\label{sec:intro}

Measuring heart rate and breathing without touching the person is one of the more mature applications of vision
outside the visible-appearance setting. Remote photoplethysmography (rPPG) recovers the pulse from subtle skin-colour
changes in ordinary video~\cite{verkruysse2008,mcduff2023}. Thermal imaging sees breathing as temperature changes at
the nostrils~\cite{ioannou2014}. Commodity Wi-Fi recovers breathing from the way a person's chest perturbs the radio
channel~\cite{liu2015,wang2016resp}, and has been used to infer emotion~\cite{zhao2016}. Progress in all three is
measured on benchmark datasets, and those datasets shape what the methods learn.

Most of those datasets record people who are sitting still and at rest~\cite{bobbia2019,stricker2014pure}. Heart
rate varies little within a session, breathing is spontaneous and unmeasured, and the conditions under which a
contactless sensor would actually be needed are absent: someone who is stressed, frightened, exercising, or
breathing unusually. Affective computing datasets have the opposite profile. They elicit emotion carefully and
record contact physiology~\cite{koelstra2012,soleymani2012,schmidt2018}, but rarely with face video, thermal video
and radio recorded together on one clock, and never with a breathing ground truth that is known by design. No
single dataset lets a model trained to read the pulse from video be tested when the same person's heart rate has
been driven up by exercise, or lets a Wi-Fi breathing estimator be scored against a breathing rate the participant
was instructed to follow.

\dataset is built to fill that gap. Its design principle is that the protocol should produce the ground truth,
rather than rely on sensors alone to supply it. Paced breathing at 6 and 20~breaths/min gives respiration targets
that bracket the normal range. Two short exercises followed by a still, two-minute recovery sweep heart rate from
rest to \ph{peak}\bpm and back, in the recovery window the physiology literature uses~\cite{cole1999}. Six emotion
film clips, each followed by a Self-Assessment Manikin rating~\cite{bradley1994}, supply dimensional and
categorical affect labels. All of this is recorded simultaneously by contactless sensors (two face cameras, a
thermal camera, Wi-Fi CSI) and contact references (ECG, PPG, EDA, EMG, audio).

Our contributions are:
\begin{itemize}
  \item \textbf{A synchronized multimodal dataset} of \ph{N} participants and \ph{H} hours, pairing lossless face
    video, thermal video and Wi-Fi CSI with ECG, PPG, EDA, EMG and audio on one host clock, with per-frame
    timestamps and a per-step quality report for every session (Sec.~\ref{sec:dataset}).
  \item \textbf{A protocol with ground truth by design}: paced breathing, an exercise-driven heart-rate sweep with
    recovery, emotion induction with self-report, and cognitive load, in one scripted 34-step session whose exact
    stimuli and trial sequences can be regenerated (Sec.~\ref{sec:protocol}).
  \item \textbf{Benchmarks and baselines} for remote heart rate, contactless respiration from three modalities,
    cross-modal Wi-Fi sensing supervised by vision, and affect recognition, with subject-independent splits
    (Sec.~\ref{sec:bench}).
  \item \textbf{Open recording software}, including the camera controls needed for rPPG-grade video (fixed
    exposure matched to mains flicker, lossless encoding, frame-accurate logging) (Sec.~\ref{sec:sync}).
\end{itemize}

\begin{figure}[t]
  \centering
  \figph{0.95\linewidth}{4.2cm}{Fig.~1 (teaser): one participant, synchronized traces over the whole session.
  Top: face / thermal / CSI thumbnails. Below: ECG heart rate, rPPG heart rate, CSI breathing band, EDA, with the
  protocol phases shaded (clips, Stroop, exercise, recovery, paced breathing 6 and 20/min, body scan).}
  \caption{\dataset records contactless sensors and contact references simultaneously while the protocol drives
  heart rate, breathing and affect through known changes. \ph{Replace with a real session, \eg participant 132.}}
  \label{fig:teaser}
\end{figure}


%% =====================================================================
%% 2_related.tex
%% =====================================================================
\section{Related Work} \label{sec:related}

\paragraph{Remote physiological sensing datasets.}
UBFC-rPPG~\cite{bobbia2019} and PURE~\cite{stricker2014pure} established camera heart-rate benchmarks with a finger
PPG reference, recorded at rest or with scripted head motion. MMSE-HR~\cite{zhang2016mmse} adds spontaneous emotion
tasks, VIPL-HR~\cite{niu2019vipl} adds scale and varied capture devices, UBFC-Phys~\cite{sabour2021ubfcphys} adds a
social stress test, and MMPD~\cite{tang2023mmpd} adds skin tone, lighting and motion diversity. SCAMPS~\cite{mcduff2022scamps}
provides synthetic faces with exact waveforms. These datasets are video-only on the sensing side, and none controls
breathing rate. Their heart rate range is mostly that of seated rest. \dataset keeps the camera-plus-contact design
and adds a controlled heart-rate sweep, paced breathing, thermal video, and radio.

\paragraph{Affective physiological datasets.}
DEAP~\cite{koelstra2012} and MAHNOB-HCI~\cite{soleymani2012} paired music-video or film stimuli with EEG and
peripheral physiology, and MAHNOB-HCI included face video. AMIGOS~\cite{miranda2018amigos} and WESAD~\cite{schmidt2018}
extended this to groups and to wearable stress detection. These datasets are the model for our stimulus and
self-report design (SAM ratings, discrete emotion labels), but they were built for affect recognition, not to test
contactless vital-sign estimation, and they do not include thermal video or Wi-Fi.

\paragraph{RF and Wi-Fi sensing.}
Wi-Fi CSI has been used for breathing and heart rhythm~\cite{liu2015,wang2016resp}, emotion~\cite{zhao2016},
gesture~\cite{zheng2019widar}, and human pose, often supervised by a camera~\cite{zhao2018rfpose,wang2019personwifi}.
MM-Fi~\cite{yang2023mmfi} pairs Wi-Fi with vision and other sensors for human sensing. Public CSI datasets with
contact physiological references recorded under emotion and exercise are, to our knowledge, not available.
\ph{Check once more for any 2025--2026 Wi-Fi vital-sign dataset before submission.}

\paragraph{Thermal physiological imaging.}
Facial thermal imaging tracks sympathetic arousal and breathing~\cite{ioannou2014}, but public thermal datasets
with synchronized contact physiology are small. \dataset adds a thermal stream under the same protocol.

\begin{table*}[t]
  \centering
  \small
  \setlength{\tabcolsep}{4pt}
  \begin{tabular}{@{}lrcccccccccc@{}}
    \toprule
    Dataset & Subj. & RGB & Thermal & Wi-Fi & ECG & PPG & EDA & Emotion & Exercise & Paced breath. & Lossless \\
    \midrule
    UBFC-rPPG~\cite{bobbia2019}          & 42\ver  & \cmark & \xmark & \xmark & \xmark & \cmark & \xmark & \xmark & \xmark & \xmark & \ph{?} \\
    PURE~\cite{stricker2014pure}         & 10\ver  & \cmark & \xmark & \xmark & \xmark & \cmark & \xmark & \xmark & \xmark & \xmark & \cmark \\
    MMSE-HR~\cite{zhang2016mmse}         & 40\ver  & \cmark & \cmark & \xmark & \xmark & \ph{?} & \cmark & \cmark & \xmark & \xmark & \ph{?} \\
    VIPL-HR~\cite{niu2019vipl}           & 107\ver & \cmark & \xmark & \xmark & \xmark & \cmark & \xmark & \xmark & \ph{?} & \xmark & \xmark \\
    UBFC-Phys~\cite{sabour2021ubfcphys}  & 56\ver  & \cmark & \xmark & \xmark & \xmark & \cmark & \cmark & \ph{stress} & \xmark & \xmark & \ph{?} \\
    MMPD~\cite{tang2023mmpd}             & 33\ver  & \cmark & \xmark & \xmark & \xmark & \cmark & \xmark & \xmark & \ph{?} & \xmark & \xmark \\
    DEAP~\cite{koelstra2012}             & 32\ver  & \ph{22 subj.} & \xmark & \xmark & \xmark & \cmark & \cmark & \cmark & \xmark & \xmark & \xmark \\
    MAHNOB-HCI~\cite{soleymani2012}      & 27\ver  & \cmark & \xmark & \xmark & \cmark & \xmark & \cmark & \cmark & \xmark & \xmark & \xmark \\
    WESAD~\cite{schmidt2018}             & 15\ver  & \xmark & \xmark & \xmark & \cmark & \cmark & \cmark & \cmark & \xmark & \xmark & -- \\
    MM-Fi~\cite{yang2023mmfi}            & 40\ver  & \cmark & \xmark & \cmark & \xmark & \xmark & \xmark & \xmark & \ph{?} & \xmark & \ph{?} \\
    \midrule
    \textbf{\dataset (ours)}             & \ph{N}  & \cmark\,$\times$2 & \cmark & \cmark & \cmark & \cmark & \cmark & \cmark & \cmark & \cmark & \cmark \\
    \bottomrule
  \end{tabular}
  \caption{Comparison with related datasets. \ver: counts and modalities taken from the original papers, to be
  verified before submission; \ph{?}: not yet checked. ``Paced breath.'' means breathing at an instructed rate that
  serves as ground truth.}
  \label{tab:compare}
\end{table*}


%% =====================================================================
%% 3_dataset.tex
%% =====================================================================
\section{The \dataset Dataset}\label{sec:dataset}

\subsection{Design goals}\label{sec:goals}
Four goals drove the design. \textbf{(G1)} Physiology must move: the protocol should drive heart rate, breathing
and arousal through large, known changes rather than hope they occur. \textbf{(G2)} Ground truth by design where
sensors cannot supply it: a visual pacer fixes breathing rate, and self-reports label every emotional stimulus.
\textbf{(G3)} Contactless and contact sensors must share one clock, with every video frame timestamped, so that a
sample in any stream can be located in every other stream. \textbf{(G4)} Recordings must be good enough for rPPG:
lossless video with exposure and white balance fixed, and never adjusted by the camera during a session.

\subsection{Acquisition system}\label{sec:system}
Table~\ref{tab:sensors} lists the sensors. All streams are captured by one acquisition program on a single Linux
host, which also receives protocol events from the stimulus display. The participant sits
\ph{50--80}\,cm from the cameras, facing a display that presents the protocol, with the Wi-Fi transmitter and
receiver on either side \ph{geometry: TX/RX positions, distance, heights, channel, bandwidth}.

\begin{table}[t]
  \centering
  \footnotesize
  \setlength{\tabcolsep}{3pt}
  \begin{tabular}{@{}l l r l@{}}
    \toprule
    Stream & Device & Rate & Format \\
    \midrule
    Face cam.\ 1 & USB webcam\ph{model} & 30\,fps & \ph{1280$\times$720}, FFV1 lossless \\
    Face cam.\ 2 & Sony a6700 (USB) & 25\,fps & 1920$\times$1080, FFV1 lossless \\
    Thermal      & FLIR UVC camera & $\approx$10\,fps & 640$\times$480, 8-bit (non-radiometric) \\
    Wi-Fi CSI    & ESP32 receiver & 100\,Hz & 57 subcarriers, I/Q \\
    ECG          & Frontier X chest strap & 125\,Hz & single lead \\
    PPG          & CMS50E finger oximeter & 60\,Hz & pleth, pulse rate, SpO$_2$ \\
    EDA          & ESP32 front end & 20\,Hz & skin conductance ($\mu$S) \\
    EMG          & ESP32 front end & 1\,kHz & raw sEMG \\
    Audio        & two lavalier mics & 48\,kHz & participant / operator channels \\
    \bottomrule
  \end{tabular}
  \caption{Sensors in \dataset. Rates are nominal; measured per-session rates are in each session's quality report.}
  \label{tab:sensors}
\end{table}

\subsection{Synchronization and recording integrity}\label{sec:sync}
\paragraph{One clock.} Every sample and every video frame is stamped on arrival with the host's monotonic clock
and wall-clock time in nanoseconds. Protocol events from the display (step onsets, clip playback start, breathing
pacer start, button presses) are sent to the host with server timestamps and converted to the same clock using a
continuously estimated offset. Clip onsets are logged when the video actually starts playing, not when the step
begins.

\paragraph{Frame accounting.} For each camera, a per-frame log records the timestamp of every frame that reached the
file, rather than every frame the camera delivered. In a pilot session the camera delivered 25\,fps while the
encoder silently dropped 42.6\% of frames; counting at the capture side would have reported a clean recording.
Encoding runs in a separate multi-threaded writer so that lossless FFV1 keeps up with the capture rate, and the
quality report checks that the frames in each video equal the rows in its log.

\paragraph{Exposure control for rPPG.} Camera auto-exposure responds to the same brightness changes the emotion
clips produce on the participant's face, which contaminates the pulse signal. The recorder lets auto-exposure
settle, then locks exposure, gain and white balance, and verifies the lock by read-back. It also sets the camera's
anti-flicker filter to the mains frequency and snaps the locked exposure to a whole number of flicker periods
(multiples of 10\,ms at 50\,Hz), so that lighting flicker integrates out of every frame. For the camera that
exposes no controls over USB, exposure, ISO and white balance are fixed on the camera body and verified from the
stability of the background brightness.

\paragraph{Contact references.} The ECG strap streams through a phone; its samples are re-stamped from the device
clock after the session, which removes the transport delay and arrival jitter of the link (\eg 4.9\,s delay and
5.2\,s jitter removed in one session). The oximeter and ECG heart rates are shown side by side during recording so
that disagreement is visible to the operator.

\paragraph{Wi-Fi null.} Before participants are seated, an empty-room CSI baseline is recorded with the same radio
configuration, so every CSI claim can be compared with a recording in which nobody breathes.

\paragraph{Quality gates.} A preflight check must pass before recording starts: every required stream present and
of good quality, camera settings locked, clock offset under 100\,ms, and enough disk for the session. During
recording, losing a required stream pauses the protocol automatically and logs why. After recording, a quality report
gives, for every step, the coverage of each required stream (minimum 95\%), pauses, and artifacts marked by the
operator.

\subsection{Protocol}\label{sec:protocol}
Each session follows the same scripted protocol of 34 steps and about \ph{37} minutes (Fig.~\ref{fig:protocol}),
presented full-screen by a stimulus display and advanced automatically except where noted.

\begin{figure*}[t]
  \centering
  \figph{0.97\textwidth}{2.6cm}{Fig.~2: protocol timeline (re-use docs/protocol/meridian-protocol-slide).
  Settle 3.8\,min | Part 1: surprise, disgust, Stroop/N-back, NASA-TLX, gate | Part 2: happy, sad, gate |
  Part 3: fear, anger, word list, read aloud | stand-to-sit, step-up | 2\,min recovery | paced 6/min, 20/min,
  box breathing | body scan, PANAS. Mark where each ground truth comes from.}
  \caption{The \dataset protocol (version 26). Colour marks what each phase contributes as ground truth: emotion
  labels (self-report), cognitive load (task performance, NASA-TLX), heart-rate sweep (exercise and recovery), and
  breathing rate (pacer).}
  \label{fig:protocol}
\end{figure*}

\paragraph{Settling and baseline.} After a sensor check and a PANAS mood questionnaire~\cite{watson1988}, the
participant rests with eyes open for 120\,s. This resting baseline is the reference for every later response and for
the return-to-baseline gates~\cite{jennings1992,taskforce1996}.

\paragraph{Emotion induction.} Participants watch six film clips of \ph{90--100}\,s, one per basic
emotion~\cite{gross1995,rottenberg2007}, in three parts ordered by the typical heart-rate response of each
emotion~\cite{siegel2018,kreibig2010}: surprise and disgust, then happiness and sadness, then fear and anger. After
each clip they rate valence, arousal and dominance on 9-point Self-Assessment Manikin scales~\cite{bradley1994},
choose the emotion they mainly felt and its intensity, and report whether they had seen the clip before. A 20\,s
neutral washout separates the two clips of a part, and Parts~1 and~2 end with a return-to-baseline gate that waits
until heart rate is within 5\bpm and skin conductance within 10\% of baseline for 15\,s. Clips are drawn from
\ph{validated film set, with citation}.

\paragraph{Cognitive load.} At the end of Part~1 the participant performs a Stroop colour-word
task~\cite{stroop1935,renaud1997} or an N-back task, alternating between groups of six participants, followed by a
NASA-TLX workload rating~\cite{hart1988}. Trial sequences are generated from a logged seed.

\paragraph{Memory and speech.} Part~3 ends with a word-list study and recognition test~\cite{snodgrass1988} and a
fixed passage read aloud, which provides read speech from every participant.

\paragraph{Heart-rate sweep and recovery.} A 60\,s stand-to-sit exercise~\cite{bohannon2019} and a 120\,s
self-paced step-up~\cite{bennett2016} raise heart rate, followed by a still, seated two-minute recovery, the window
in which heart-rate recovery is conventionally measured~\cite{cole1999}. Participants who report a heart, breathing
or mobility condition in the safety screening skip both exercises, and the skip is logged.

\paragraph{Paced breathing.} A visual pacer (a rising and falling marker, with no audio) sets breathing at
6\,breaths/min for 60\,s and 20\,breaths/min for 42\,s, followed by 64\,s of box breathing. Each condition starts when
the participant presses Start, and the export records that press as its onset. The two paced rates bracket the normal
range, from resonance-frequency breathing~\cite{lehrer2014} to fast breathing. They are the respiration ground
truth for the contactless benchmarks. Box breathing~\cite{balban2023} is a relaxation condition and is not used for
scoring.

\paragraph{Relaxation.} A 270\,s guided body scan~\cite{kabatzinn1990} and a short final rest precede a second
PANAS, so the session ends in a deliberately low-arousal state.

\subsection{Participants and ethics}\label{sec:participants}
\ph{N} adults (\ph{age range, mean$\pm$sd; sex; Fitzpatrick skin type distribution; eyewear, facial hair}) took part
after written informed consent, under a protocol approved by \ph{institutional ethics board, approval number, in
the camera-ready version only}. Participants chose, at consent, which parts of their data may be shared (\eg face
video, audio), and the release honours these choices per file. Participants were screened for conditions
contraindicating exercise, and no names are stored with the data.

\subsection{Labels and release format}\label{sec:release}
Each session is released as raw streams, a per-frame timestamp log for every camera, a protocol event table, and
an analysis bundle in BIDS layout~\cite{gorgolewski2016bids}. The bundle contains per-step windows, per-clip
self-reports (SAM valence, arousal, dominance, discrete emotion, intensity, familiarity), questionnaires (PANAS
start and end, NASA-TLX, TIPI personality~\cite{gosling2003}), task trial tables (Stroop or N-back, word list),
derived signals (ECG R-peaks, per-step heart rate and heart-rate variability, skin conductance level), operator
artifact markers (talking, movement, cough, operator in view), and the quality report. Every file is checksummed
in a manifest. \ph{Licence, access procedure (data use agreement for identifiable video), hosting.}


%% =====================================================================
%% 4_statistics.tex
%% =====================================================================
\section{Dataset Statistics and Quality}\label{sec:stats}

\paragraph{Scale.} \ph{N} sessions, \ph{H} hours of synchronized recording, \ph{F} million lossless face frames,
\ph{size} TB of raw data (one session is about 180\,GB, dominated by the two lossless cameras).

\paragraph{Physiological range.} Fig.~\ref{fig:stats}a shows the ECG heart-rate distribution per protocol phase.
\req{HR per phase: rest, clips, Stroop/N-back, exercise peak, recovery, paced breathing, body scan. Headline:
median and 95th percentile HR during exercise vs rest, and the range covered overall.} The recovery segment gives
\ph{mean HR drop in the first 60\,s} per participant.

\paragraph{Affect labels.} \req{SAM valence--arousal scatter per clip (Fig.~\ref{fig:stats}b); percentage of clips
whose reported emotion matches the target; PANAS change from start to end.}

\paragraph{Recording quality.} \req{Per-stream coverage across all steps; frame drops per camera (expected
$\approx$0\%); number of automatic pauses per session and their causes; ECG--PPG heart-rate agreement.}

\paragraph{Known limitations of individual streams.} We report these as measured rather than hide them. The
surface EMG is dominated by 50\,Hz mains pickup (85--97\% of in-band power in a bench test), so we release it raw
with that caveat. The thermal camera provides 8-bit video without radiometric temperature. There is no respiration
belt: breathing ground truth comes from the pacer, and spontaneous breathing outside the paced steps has no contact
reference. Finger-PPG coverage has gaps where the oximeter's USB link dropped, and these are flagged in the quality
report. There is no photodiode on the display, so display latency is not measured per session.
\ph{Update after the full collection.}

\begin{figure}[t]
  \centering
  \figph{0.95\linewidth}{3.6cm}{Fig.~3: (a) ECG heart rate per protocol phase (violin); (b) SAM valence--arousal
  per clip, coloured by target emotion.}
  \caption{\dataset statistics. \ph{Fill after collection.}}
  \label{fig:stats}
\end{figure}


%% =====================================================================
%% 5_benchmarks.tex
%% =====================================================================
\section{Benchmarks}\label{sec:bench}

All benchmarks use subject-independent splits: \ph{k}-fold cross-validation over participants, with fold
assignments released. Signals are windowed per protocol step using the logged step onsets.

\subsection{T1: Remote heart rate}
\textbf{Task.} Estimate heart rate from face video in \ph{10}\,s windows. \textbf{Reference.} ECG R-peaks.
\textbf{Metrics.} MAE, RMSE (bpm), MAPE, Pearson $r$, reported overall and per phase (rest, emotion, exercise,
recovery, paced breathing). \textbf{Baselines.} Unsupervised GREEN~\cite{verkruysse2008},
CHROM~\cite{dehaan2013chrom} and POS~\cite{wang2017pos}; supervised PhysNet~\cite{yu2019physnet},
TS-CAN~\cite{liu2020tscan} and PhysFormer~\cite{yu2022physformer}, trained on \dataset and, for cross-dataset
transfer, on UBFC-rPPG and PURE, using rPPG-Toolbox~\cite{liu2023toolbox}.
\textbf{Question the split answers.} How much of a method's accuracy at rest survives the exercise sweep?

\begin{table}[t]
  \centering
  \footnotesize
  \setlength{\tabcolsep}{3pt}
  \begin{tabular}{@{}l ccc ccc@{}}
    \toprule
    & \multicolumn{3}{c}{Rest + clips} & \multicolumn{3}{c}{Exercise + recovery} \\
    \cmidrule(lr){2-4}\cmidrule(lr){5-7}
    Method & MAE$\downarrow$ & RMSE$\downarrow$ & $r\uparrow$ & MAE$\downarrow$ & RMSE$\downarrow$ & $r\uparrow$ \\
    \midrule
    GREEN~\cite{verkruysse2008} & \ph{} & \ph{} & \ph{} & \ph{} & \ph{} & \ph{} \\
    CHROM~\cite{dehaan2013chrom} & \ph{} & \ph{} & \ph{} & \ph{} & \ph{} & \ph{} \\
    POS~\cite{wang2017pos}      & \ph{} & \ph{} & \ph{} & \ph{} & \ph{} & \ph{} \\
    PhysNet~\cite{yu2019physnet} & \ph{} & \ph{} & \ph{} & \ph{} & \ph{} & \ph{} \\
    TS-CAN~\cite{liu2020tscan}  & \ph{} & \ph{} & \ph{} & \ph{} & \ph{} & \ph{} \\
    PhysFormer~\cite{yu2022physformer} & \ph{} & \ph{} & \ph{} & \ph{} & \ph{} & \ph{} \\
    \bottomrule
  \end{tabular}
  \caption{T1 remote heart rate on face camera~1, subject-independent. \req{run rPPG-Toolbox baselines}.}
  \label{tab:t1}
\end{table}

\subsection{T2: Contactless respiration}
\textbf{Task.} Estimate breathing rate during the two paced steps. \textbf{Reference.} The pacer rate (6 and
20\brpm); a participant's step is excluded if \ph{adherence criterion, \eg the dominant chest-motion frequency in
video is more than 1\brpm from the pacer}. \textbf{Modalities and baselines.} Video (optical flow on the upper
body~\cite{mcduff2023}), thermal (nostril-region temperature), Wi-Fi CSI (spectral peak of the principal component
of subcarrier amplitudes after removing the empty-room baseline~\cite{wang2016resp,liu2015}), and fusion.
\textbf{Metric.} MAE in breaths/min per rate. \req{T2 table: modality $\times$ \{6, 20\}\brpm, MAE and success
rate within 1\brpm.}

\subsection{T3: Cross-modal Wi-Fi sensing supervised by vision}
Face video is a rich but privacy-sensitive signal. Wi-Fi is private but hard to label. Because \dataset records
both on one clock, the camera can supervise the radio. We provide a simple baseline, \method. A CSI window $\csi$ is
encoded into $\zwifi$, and a frozen visual respiration model gives an embedding $\zvis$ and a breathing waveform
$\svis$ from the synchronized video. \method is trained with
\begin{equation}
  \mathcal{L} = \Ltask + \lambda_{a}\,\Lalign\!\left(\zwifi,\zvis\right) + \lambda_{w}\,\Lwave\!\left(\hat{\mathbf{s}}^{\mathrm{wifi}},\svis\right),
  \label{eq:respwave}
\end{equation}
where $\Ltask$ is the rate loss on labelled paced steps, $\Lalign$ aligns the Wi-Fi and visual embeddings
\ph{contrastive / cosine}, and $\Lwave$ matches the predicted Wi-Fi breathing waveform to the visual one
\ph{negative Pearson}. At test time only Wi-Fi is used. \req{T3: CSI-only vs \method, trained with visual
supervision on unpaced steps and evaluated on paced steps, MAE in brpm; ablations of $\lambda_a$, $\lambda_w$.}

\subsection{T4: Affect recognition}
\textbf{Task.} Classify each clip window as high/low valence and arousal (split at the scale midpoint) and as one of
six target emotions. \textbf{Labels.} Participant's own SAM ratings and reported emotion. \textbf{Baselines.}
Per-modality features (heart rate and heart-rate variability from ECG, skin conductance level and responses, facial
expression features from face video, thermal nose and forehead temperature, CSI statistics) with gradient-boosted
trees, and late fusion. \textbf{Metrics.} Accuracy and macro-F1, leave-participants-out. \req{T4 table: modality
$\times$ \{valence, arousal, 6-class\}.}


%% =====================================================================
%% 6_limitations.tex
%% =====================================================================
\section{Limitations and Ethical Considerations}\label{sec:limits}

\textbf{Fixed stimulus order.} All participants see the clips in the same order, graded by expected heart-rate
response, so session-level drift and habituation are confounded with emotion. Gates and washouts limit carry-over
but do not remove the confound. \textbf{Breathing after exercise.} The paced steps follow the exercise recovery, so
they are valid as a breathing-rate reference but heart-rate variability measured there is not a resting value.
\textbf{References.} There is no respiration belt or display photodiode (Sec.~\ref{sec:stats}).
\textbf{Population.} \ph{N} participants recruited from \ph{population}; skin-type and age coverage are reported in
Sec.~\ref{sec:participants} and limit the claims we can make about fairness of rPPG across skin tones.
\textbf{Privacy.} Face video and voice are identifiable. They are released only under a data use agreement and only
for participants who consented to sharing them. Wi-Fi, thermal and physiological streams are released more
openly. Users must not attempt re-identification.


%% =====================================================================
%% 7_conclusion.tex
%% =====================================================================
\section{Conclusion}\label{sec:conclusion}
\dataset pairs face video, thermal video and Wi-Fi CSI with contact physiology, recorded on one clock while a
scripted protocol drives emotion, cognitive load, heart rate and breathing through known changes. Its benchmarks
test contactless sensing where it is hardest and most useful, beyond the calm, seated recordings current
benchmarks are built on. \ph{One sentence on what the baselines revealed.}


{
    \small
    \bibliographystyle{ieeenat_fullname}
    \bibliography{main}
}

\end{document}
